% Página tras la portada: colaboradores, estructura del repositorio y notación.
\section*{Sobre estos apuntes}

Estos apuntes se mantienen en el repositorio
\href{https://github.com/jmnu4245/Topologia-Algebraica}{\texttt{github.com/jmnu4245/Topologia-Algebraica}}.
Cada vez que se actualiza la rama \texttt{main} se compila el PDF y se publica en la sección
\href{https://github.com/jmnu4245/Topologia-Algebraica/releases}{\emph{Releases}}. Las erratas y sugerencias se recogen como \emph{issues} del repositorio.

\subsection*{Coautores}
Personas que han contribuido abriendo \emph{issues}:
\begin{itemize}
    \item \textbf{DavidP2004}
    \begin{itemize}
        \item Corrección de la definición de conjunto convexo (\#1).
        \item Sugerencia de demostración del inverso homotópico de caminos (\#3).
        \item Error en el enunciado de la unicidad de levantamientos (\#4).
        \item Demostración de que $S^1$ no es contráctil (\#5).
        \item Error en un ejercicio (\#6).
    \end{itemize}
    \item \textbf{T1N8T}
    \begin{itemize}
        \item Errata $[\gamma] \in C_{x_1}(X)$ en la demostración del isomorfismo de cambio de base (\#7).
        \item Errata $C_{x_0}$ por $c_{x_0}$ en la demostración de que un estrellado tiene $\pi_1$ trivial (\#8).
        \item Propuesta de marcar los resultados demostrados en clase con $\odot$ (\#10).
    \end{itemize}
\end{itemize}

\subsection*{Contenido del repositorio}
\begin{itemize}
    \item \textbf{Raíz:} \texttt{main.tex} es el archivo principal y \texttt{preambulo.tex} contiene los paquetes, los estilos de las cajas y las macros.
    \item \textbf{\texttt{teoria/}}: el contenido de los apuntes.
    \begin{itemize}
        \item \texttt{presentacion.tex} es esta página e \texttt{intro.tex} la información de la asignatura.
        \item Hay un archivo por tema:
        \begin{enumerate}
            \item \texttt{ResultadosPreliminares.tex}
            \item \texttt{EquivalenciaHomotopica.tex}
            \item \texttt{ElGrupoFundamental.tex}
            \item \texttt{EspaciosSimplementeConexosyGrupoFundamental.tex}
            \item \texttt{AplicacionesTeoriaHomotopia.tex}
            \item \texttt{TeoremaSeifertVanKampen.tex}
        \end{enumerate}
        \item \texttt{ResumenHomotopias.tex} es la hoja resumen de homotopías, lazos y grado, después del último tema.
    \end{itemize}
    \item \textbf{\texttt{ejercicios/}}: un archivo por boletín (\texttt{tema1.tex}, \dots, \texttt{tema5.tex}), \texttt{ejercicios.tex}, que los reúne, y \texttt{examenes.tex}.
    \item \textbf{\texttt{papel/}}: fotografías de los apuntes manuscritos de clase, a partir de los cuales se ha redactado este documento. Hay una carpeta por clase, con nombre \texttt{aaaa-mm-dd\_claseN}.
    \item \textbf{\texttt{practicas/}}: apuntes de la parte de prácticas (\texttt{main.tex}), los cuadernos de cada práctica con Singular (\texttt{practN.ipynb}) y las fotografías de las prácticas.
    \item \textbf{\texttt{web/}}: versión web, publicada en \href{https://jmnu4245.github.io/Topologia-Algebraica/}{\texttt{jmnu4245.github.io/Topologia-Algebraica}}: los apuntes en HTML y un mapa interactivo de resultados. \texttt{construir.sh} la genera; GitHub Actions la publica en cada cambio.
    \item \textbf{\texttt{mapa\_teoremas/}}: versión antigua del mapa de dependencias, en PDF.
    \item \textbf{\texttt{scripts/}}: \texttt{exportar.sh} genera en \texttt{exportados/} un PDF por tema (teoría y ejercicios por separado) y otro de los exámenes; \texttt{installsingular.sh} instala Singular, usado en las prácticas.
    \item \textbf{\texttt{.github/workflows/}}: compilación automática del PDF y de la web.
\end{itemize}

\subsection*{Notación}
\begin{itemize}
    \item \textbf{Fechas.} Un encabezado \textbf{dd/mm/aaaa (Clase $n$)} indica que el contenido que sigue se vio en la clase $n$, impartida en esa fecha.
    \item \textbf{Referencias a las clases.} Los títulos de las cajas terminan con una referencia entre corchetes:
    \begin{itemize}
        \item \textbf{[C$n$.$m$]} es el punto $m$ de la clase $n$. Por ejemplo, [C8.11] es el punto 11 de la clase 8.
        \item \textbf{[C$n$.$m$–$n$.$m'$]} abarca varios puntos de esa clase.
        \item \textbf{[C$n$]} indica que el resultado es de la clase $n$, sin un punto concreto.
    \end{itemize}
    No hay que confundir esta referencia con la numeración propia del documento: el \emph{Teorema 4.3} es el tercer teorema de la sección 4.
    \item \textbf{Resultados demostrados en clase.} Los resultados cuya demostración se hizo en clase llevan en el título una marca $\odot$ con una letra:
    \begin{itemize}
        \item \fcolorbox{white}{red!80!black}{\textcolor{white}{\textbf{$\odot$\,T}}} para un teorema;
        \item \fcolorbox{white}{red!80!black}{\textcolor{white}{\textbf{$\odot$\,C}}} para un corolario;
        \item $\odot$\,P para una proposición;
        \item $\odot$\,L para un lema.
    \end{itemize}
    Los teoremas y corolarios van destacados en rojo.
    \item \textbf{Ejercicios propuestos en clase.} Los resultados que el profesor dejó como ejercicio en clase (``Ejercicio:'' en los contenidos de las clases) llevan la marca \marcaE. Están demostrados en su sitio, dentro de cada tema, porque hacen falta para seguir la teoría.
    \item \textbf{Ejercicios de los boletines.} Están todos al final, en la sección \emph{Ejercicios de los boletines}, en cajas verdes numeradas como en el boletín (T4.2 es el ejercicio 2 del Tema 4). Si la solución ya está en los apuntes, solo se da el enlace; si no, se indica de dónde sale (clase, examen o solución propia). La sección \emph{Exámenes de 2025-26} dice dónde está la solución de cada pregunta de examen.
\end{itemize}
